\qhead{2. Making Pricing Decisions, Value and Prices. What are your product's reference and differential values? What are your price floor and price ceiling? Explain your answer.}

To judge what the Qin Plus DM-i is worth I use the exchange-value idea, where the most a buyer will rationally pay is the reference value plus the differentiation value \cite{nagle2018}.

The reference value is the price of whatever the customer would otherwise buy, and for an ordinary Chinese family that is not another electric car but a petrol one, a Volkswagen Lavida or a Toyota Corolla somewhere around \euro13{,}000 to \euro17{,}000. That figure is the anchor. Sitting on top of it is the differentiation value, which comes in two kinds. The monetary part is real money the owner keeps or avoids spending: fuel savings that can reach about \euro1{,}240 a year, the exemption from the new-energy-vehicle tax, and lower maintenance over the life of the car. The psychological part is harder to put a number on but still counts, and it covers the modern technology, the driver-assistance system BYD includes for nothing, and the feeling of doing something for the environment. Because both parts are positive, the exchange value ends up comfortably above the price of the petrol car.

That gives the two limits on price. The floor is not full accounting cost but the extra cost of building one more car, and vertical integration together with the Blade Battery, which is roughly ten euros per kilowatt-hour cheaper than the older nickel-cobalt cells \cite{motleyfool2025}, pushes that floor unusually low. The ceiling is the exchange value itself, because once the price climbs past it the buyer keeps more by going back to petrol. The space between the two is the consumer surplus, the gap between what the car is worth to the buyer and what they actually hand over, and BYD deliberately leaves a wide one.

What it does is price the Qin below the petrol car it is being measured against while giving more value, which is exactly what ``electricity cheaper than oil'' claims out loud. The point worth making is that this is not cost-driven pricing in disguise. The low cost is what lets BYD survive at these prices, but the price itself is set against the petrol alternative. Had the company simply marked up its low cost, it would have priced higher and missed the very buyers it set out to win over. The same logic raises the ceiling in Europe, where the petrol cars used as the reference are dearer and the buyers wealthier, which is why the identical model can sit much higher up its price window abroad.
